\chapter{Introduction}
\label{sec:introduction}
Automatic speech recognition (ASR) has improved substantially in recent years due to advances in deep learning, large-scale training data, and computational resources. However, performance remains much less reliable in realistic conversational settings, particularly when multiple speakers overlap, and background noise is present. To develop and evaluate models under these conditions, robust, realistic datasets are needed. The AMI meeting corpus \cite{ami-corpus} and the CHIME-6 dataset \cite{chime6} provide important examples of real-world multi-speaker datasets with realistic, detailed annotations.

Two major factors that degrade ASR performance are overlapping speech and environmental noise. Both are common in natural human conversation \citep{chen2020continuousspeechseparationdataset, Conversation_in_small_groups}. Overlapping speech is the time interval during which two or more speakers speak simultaneously. At the same time, background noise is any unintended sound that does not contribute to the conversation. Overlapping speech violates the single-active-speaker assumption underlying many conventional ASR systems. Natural conversations often involve single-speaker segments and partially overlapping regions \cite{chen2020continuousspeechseparationdataset}, where standard metrics such as word error rate (WER) might fail to reflect the accuracy \citep {ellis2026werunawareassessingasr, hughes2025problemwithwer}. AS a result, overlap-aware and speaker-aware measures such as concatenated minimum-permutation WER (cpWER), Optimal Reference Combination WER (ORC-WER), and diarisation invariant cpWER (DI-cpWER) are proposed \citep{ Word_Error_Rate_Definitions}.

Due to the difficulties and cost of collecting and annotating real-world multi-speaker datasets, synthetic mixtures have emerged as an important alternative. Existing datasets have improved realism by adding ambient noise, reverberation and increasing the diversity of recording conditions and scenarios. However, most were fixed rather than flexible for downstream training to systematically produce more data and varying overlap, noise, duration and speaker count.

Despite this progress, a practical gap remains.  Real meeting conversations preserve natural patterns, but this same naturalness poses lexical and acoustic challenges, limiting control over the degradation factors during training and evaluation. On the other hand, existing synthetic datasets are easier to manipulate, but they are typically built with fixed or partially controlled parameters. As a result, it remains difficult to isolate how factors such as overlap ratio, signal-to-noise ratio, mixture duration and speaker count individually and jointly affect ASR performance.

This dissertation addresses the gap by introducing a synthetic mixture generation framework that takes clean speaker recordings, transcription and noise recording as input and outputs mixtures with a user-specified overlap ratio, signal-to-noise ratio, audio length and speaker count, along with timestamped transcripts. Utilising this framework, this study investigates three main questions: how ASR performance changes with overlap and noise, how transferable is the trend observed in synthetic data to real-world, and how publicly available ASR models compare under different conditions.

This framework provides more transparency and information in model evaluation by separating different conditions by degradation factors. In the case where robustness changes with degradation conditions, single-point evaluation becomes underrepresentative. Hence, instead of simply reporting under a single dataset, this framework provides an opportunity to analyse how different models behave in a range of different conditions, identifying the strengths and weaknesses of different models.

The structure of the dissertation is as follows: Chapter 2 reviews the related work on ASR robustness, multi-speaker datasets and evaluation metrics. Chapter 3 describes architectural design choicesused in the experiments. Chapter 4 details the experimental design, including the generation of synthetic mixtures and the evaluation on real-world data. Chapter 5 outlines and reports the results of experiments. Chapter 6 discusses the implications of the findings, limitations of the study and directions for future work. Finally, Chapter 7 concludes the dissertation by summarising the key contributions and insights gained from the research.